% macros.tex: preamble and macros for the OpenBench Functional Test Case Descriptions
% (layout copied from docs/design/macros.tex on 2026-11-05 and reduced to what this document uses).
%
% Test cases
%   \tcarea{DRF}                                  sets the area code of the test cases that follow and restarts their numbering
%   \begin{testcase}{key}{Title}{Req.Id, Req.Id}  one test case, numbered TC-<area>-nn, label tc:<key>
%     \tcfield{Label}{text}                       a field row (Actor, Priority, Preconditions, Test data, Pass criteria ...)
%     \tcsteps                                    the header row of the step table; restarts step numbering
%     \tcstep{Action}{Input}{Expected result}     one numbered step
%   \end{testcase}
%   \TC{key}                                      hyperlinked reference to a test case (prints TC-DRF-03)
%   \req{Req.Id}                                  an SRS requirement identifier, linked to its row in the traceability matrix
%   \tbd{text}                                    a visible placeholder; none may remain in the submitted edition
%
% The traceability matrix (build/traceability.tex) is generated by tools/gen_trace.py from the \begin{testcase} lines of
% sections/cases/*.tex and the \fr lines of the SRS short edition, so the requirement list of each test case is the only
% place its coverage is recorded.

\usepackage[letterpaper,margin=1in]{geometry}
\usepackage[T1]{fontenc}
\usepackage{lmodern}
\usepackage{microtype}
\usepackage{parskip}
\usepackage{enumitem}
\usepackage{booktabs}
\usepackage{longtable}[=v4.13]   % v4.24 (TeX Live 2025) raises "Infinite glue shrinkage" on every multi-page table; v4.13 does not
\usepackage{array}
\usepackage{tabularx}
\usepackage{multirow}
\usepackage{makecell}
\usepackage[table]{xcolor}
\usepackage{graphicx}
\usepackage{caption}
\usepackage{fancyhdr}
\usepackage{lastpage}
\usepackage{needspace}
\usepackage{etoolbox}
\usepackage[colorlinks=true,linkcolor=blue!50!black,urlcolor=blue!50!black,citecolor=blue!50!black]{hyperref}

% ---- Document metadata -------------------------------------------------------
\newcommand{\TCDversion}{0.2}
\newcommand{\TCDdate}{2026-11-07}
\newcommand{\SRSversion}{0.4}       % the requirements this document tests
\newcommand{\SDDversion}{1.1}       % the design the product follows
\newcommand{\SRSproduct}{OpenBench}

% ---- Layout ------------------------------------------------------------------
\setcounter{secnumdepth}{3}
\setcounter{tocdepth}{2}
\setlist{itemsep=1pt,topsep=2pt,parsep=0pt}
\setlist[enumerate]{leftmargin=*}
\setlist[itemize]{leftmargin=*}
\captionsetup{font=small,labelfont=bf}
\graphicspath{{build/figures/}}
\newcolumntype{L}[1]{>{\raggedright\arraybackslash}p{#1}}
\newcolumntype{Y}{>{\raggedright\arraybackslash}X}

\pagestyle{fancy}
\fancyhf{}
\lhead{\small \SRSproduct{} Functional Test Case Descriptions}
\rhead{\small Version \TCDversion}
\cfoot{\small Page \thepage\ of \pageref*{LastPage}}
\renewcommand{\headrulewidth}{0.4pt}
\fancypagestyle{plain}{\fancyhf{}\cfoot{\small Page \thepage\ of \pageref*{LastPage}}\renewcommand{\headrulewidth}{0pt}}

\newlength{\reqgap}\setlength{\reqgap}{2pt}

% ---- Small helpers -------------------------------------------------------------
\newcommand{\code}[1]{\texttt{#1}}
\newcommand{\term}[1]{\emph{#1}}
\newcommand{\tbd}[1]{\textcolor{red!70!black}{\textbf{[TBD]}~#1}}
\newcommand{\stretchtag}{\textcolor{black!60}{[Stretch]}}
\newcommand{\req}[1]{\hyperref[req:#1]{\textbf{\small #1}}}

% ---- Glossary entries (emitted by ../srs/tools/gen_glossary.py --compact) ------
\newcommand{\glossterm}[3]{%
  \par\noindent\hangindent=1.5em\hangafter=1%
  \textbf{#1}\quad #2%
  \if\relax\detokenize{#3}\relax\else\ {\small\textit{Avoid:} #3.}\fi\par\vspace{\reqgap}}
\newcommand{\glossgroup}[1]{\par\medskip\needspace{5\baselineskip}\noindent\textbf{#1}\par\nobreak\smallskip}

% ---- Test cases ------------------------------------------------------------------
% One longtable per test case: a title bar, field rows (label in the left column, text across the other three),
% then the numbered steps (Action, Input, Expected result). The table may break across pages; the title repeats.
\makeatletter
\newcounter{tc}
\newcounter{tcstep}
\def\tc@area{XXX}
\newcommand{\tcarea}[1]{\gdef\tc@area{#1}\setcounter{tc}{0}}
\renewcommand{\thetc}{TC-\tc@area-\two@digits{\value{tc}}}
\newcommand{\TC}[1]{\hyperref[tc:#1]{\ref*{tc:#1}}}

\newlength{\tc@w}        % text width of the Action, Input, and Expected result columns together
\newlength{\tc@lab}      % label / step-number column
\newlength{\tc@rest}     % a multicolumn over those three columns (their widths plus the padding and rules between them)
\newlength{\tc@full}     % a multicolumn over all four columns
\newlength{\tc@a}\newlength{\tc@b}\newlength{\tc@c}
\setlength{\tc@lab}{2.3cm}
\colorlet{tcbar}{black!10}

\def\tc@sep{}
\newcommand{\tc@onereq}[1]{\tc@sep\req{#1}\gdef\tc@sep{, }}
\newcommand{\tc@reqs}[1]{\gdef\tc@sep{}\forcsvlist{\tc@onereq}{#1}}

\newenvironment{testcase}[3]{%
  \refstepcounter{tc}\label{tc:#1}%
  \setcounter{tcstep}{0}%
  \setlength{\tc@w}{\dimexpr\linewidth-8\tabcolsep-5\arrayrulewidth-\tc@lab\relax}%
  \setlength{\tc@a}{0.31\tc@w}\setlength{\tc@b}{0.26\tc@w}\setlength{\tc@c}{\dimexpr\tc@w-\tc@a-\tc@b\relax}%
  \setlength{\tc@rest}{\dimexpr\tc@w+4\tabcolsep+2\arrayrulewidth\relax}%
  \setlength{\tc@full}{\dimexpr\linewidth-2\tabcolsep-2\arrayrulewidth\relax}%
  \needspace{24\baselineskip}%   the title, the fields, and the first step start on one page
  \phantomsection\addcontentsline{toc}{subsubsection}{\thetc\quad #2}%
  \small
  \setlength{\LTpre}{4pt}\setlength{\LTpost}{6pt}%
  \begin{longtable}{|L{\tc@lab}|L{\tc@a}|L{\tc@b}|L{\tc@c}|}
  \hline
  \rowcolor{tcbar}\multicolumn{4}{|L{\tc@full}|}{\normalsize\textbf{\thetc\quad #2}} \\* \hline
  \endfirsthead
  \hline
  \rowcolor{tcbar}\multicolumn{4}{|L{\tc@full}|}{\normalsize\textbf{\thetc\quad #2} \textit{\small(continued)}} \\ \hline
  \endhead
  \multicolumn{4}{r}{\footnotesize\itshape continues on next page} \\
  \endfoot
  \endlastfoot
  \textbf{Requirements} & \multicolumn{3}{L{\tc@rest}|}{\tc@reqs{#3}} \\* \hline
}{%
  \end{longtable}\normalsize\vspace{2pt}}

% Field rows and the step header end with \\* (no page break after them), so a test case never breaks inside its
% header block; the table breaks only between steps and before the closing fields.
\newcommand{\tcfield}[2]{\textbf{#1} & \multicolumn{3}{L{\tc@rest}|}{#2} \\* \hline}
\newcommand{\tcsteps}{%   \rowcolor must come first in the row, so the step counter restarts inside the cell
  \rowcolor{tcbar}\textbf{Step}\setcounter{tcstep}{0} & \textbf{Action} & \textbf{Input} & \textbf{Expected result} \\* \hline}
\newcommand{\tcstep}[3]{\stepcounter{tcstep}\arabic{tcstep} & #1 & #2 & #3 \\ \hline}

% Short edition (written by tools/gen_short.py): a test case kept in the full edition only still takes its number, so
% identifiers are the same in both editions; \tcomittedlist ends the area with those cases, where \TC{key} links to them.
\def\tc@omitted{}
\newcommand{\tcomitted}[2]{\stepcounter{tc}\xappto\tc@omitted{\noexpand\tc@omitentry{#1}{\thetc}{\unexpanded{#2}}}}
\newcommand{\tc@omitentry}[3]{\tc@sep\phantomsection\def\@currentlabel{#2}\label{tc:#1}\textbf{#2}~#3\def\tc@sep{; }}
\newcommand{\tcomittedlist}{%
  \ifdefempty{\tc@omitted}{}{%
    \par\smallskip\noindent{\textit{In the full edition only:} \def\tc@sep{}\tc@omitted.\par}%
    \gdef\tc@omitted{}}}
\makeatother
